\begin{viuabstract}
Este Trabajo Fin de Master estudia la coordinacion distribuida de robots moviles autonomos para transportar cargas heterogeneas en entornos intralogisticos. La tesis formula el problema como una cadena fisico-economica: una carga no demanda solo robots, sino capacidad efectiva, contactos utiles, fuerza y torque, movimiento seguro, comunicacion local y recuperacion ante eventos. Sobre esa lectura se evaluan familias clasicas, centralizadas, de mercado, poblacionales, primal-dual, aprendidas y de control.

La propuesta se articula alrededor de una senal comun de deficit y de la arquitectura Smith-QR, que combina dinamicas poblacionales, memoria local, estimacion de ocupacion, cierre entero de coaliciones y verificacion de capacidad en espacio \emph{wrench}. La validacion se organiza en ocho estudios SP1--SP8: reclutamiento, capacidad efectiva, factibilidad \emph{wrench}, movimiento, transporte cooperativo, robustez, comunicacion degradada y escalabilidad. Las campanas canonicas suman mas de 160.000 evaluaciones/comprobaciones, con mundos compartidos, contrastes pareados, correccion de Holm, auditorias sin fallos y videos de inspeccion cualitativa.

Los resultados muestran que cardinalidad, capacidad escalar y factibilidad fisica no son equivalentes; que los enfoques aprendidos son comparadores utiles pero no sustituyen la trazabilidad del modelo; y que las familias distribuidas o jerarquicas conservan valor cuando las referencias centralizadas pierden practicidad a escala. El alcance se mantiene acotado a simulacion reproducible e integracion visual: no se declara validacion en hardware ni despliegue industrial. Esta delimitacion permite una contribucion ambiciosa, falsable y defendible.

\keywords{AMR, transporte cooperativo, capacidad efectiva, wrench, teoria de juegos.}
\end{viuabstract}
